\scanpage{59}
\mspage{50}
% archive sheet number: 59
Le point délicat ici\unsure{,|;} est apparemment dans la notion de passage en \unsure{quotient}, \deleted{\unsure{plus cartes} ou contraction} qui nous fait sortir de la notion de ``homomorphisme \unsure{pousse-bouton|presse-bouton}.

Dans tous ces exemples, une multistrate est \unsure{transformée|transformé} en \unsure{item} (cas de \inserted{$f^*$ pour} $f$ surjectif, ou de $i_*$), ou bien en item ou $\unsure{\varnothing|\phi}$ (cas de $f^*$ en général). Il y a des cas qui s'apparentent au ``plongement'' d'un atlas dans un autre, où multistrate n'est \inserted{pourtant} pas \unsure{transformé|transformée} en item :

\uline{Exemple 2} Soit $L,L'$ deux ens, $L'\ne\varnothing$ considérons $\unsure{F''|F'}\in\mathrm{Figures}(L')$, \inserted{$F'\ne\varnothing$,} et
\diagram{\begin{tikzcd}[ampersand replacement=\&,column sep=large,row sep=large]
\mathrm{Fig}(L) \arrow[r,dash] \& \mathrm{Fig}(L\times L') \\
F \arrow[r,mapsto] \& F\times F'=\mathrm{pr}_1^*(F).\mathrm{pr}_2^*(F')
\end{tikzcd}}
C'est le composé d'un $\mathrm{pr}_1^*\deleted{\unsure{\circ}}$ (cas déjà regardé, $\mathrm{pr}_1$ surjectif, et d'une application interne à $\mathrm{Fig}(L\times L')$, savoir la
\[
\Phi\longmapsto\Phi.F\qquad\text{(où }F=\mathrm{pr}_2^*(F')\text{)}.
\]
C'est là un autre exemple d'une application qui (si $\operatorname{supp}F$ est $\unsure{\ne V|\ne\varnothing}$) a tendance à être un vrai plongement \inserted{\unsure{sauf|sans} \unsure{pour} figures (multistrates) \unsure{vides}}
